\section{Discussion}\label{sec:discussion}
\paragraph{What the bound reveals about the data.} The exact-topology bound is a property of a field and a direction, not of a layout, and it is informative on its own. A small $\nu^*$, as on Wildfires and Storms, says that the hierarchy of the field and its geography largely agree: sibling subtrees occupy separate stretches of the axis, and any good order can draw features near their locations. A large $\nu^*$, as on hourly PM$_{2.5}$, says that nested maxima interleave in space; there, a map that keeps the tree must displace features, and the displacement measures how far the merge structure of the field departs from its spatial arrangement. Because the bound requires no layout, it can be computed for several candidate directions before a map is drawn and used to choose the axis along which hierarchy and space conflict least.

\paragraph{Forced error versus layout choice.} The gap between existing maps and the bound is the most practical finding. Most of the order error of TMTM and ST-MTM on our real-data sequences is not forced by the tree; it follows from ordering leaves without a spatial reference. Any merge tree map that keeps the tree, including future ones, can be measured against the same bound and judged on how much of its displacement it chooses.

\paragraph{Spending topological fidelity deliberately.} A merge budget changes merge levels by a declared amount; it never removes a feature, alters an extremum value or moves the data. Persistence simplification at a comparable threshold would delete features below it instead. The per-frame records show where the budget was spent, and the share of frames that use it distinguishes sequences whose conflicts are rare and deep, such as Storms, from sequences in which nearly every frame is entangled, such as PM$_{2.5}$ and the volumes. Marking the changed barriers in the map can keep this trade visible to the reader.

\paragraph{One axis, several readings.} A fixed spatial axis and a geometry-oriented embedding answer different questions. AnchorMTM supports reading position and motion along a chosen direction, ST-MTM preserves local neighborhoods better, and the central-tree variant of Fig.~\ref{fig:cylinder-centered} exposes regime changes. These readings complement one another, and a view that links them, or small multiples along several directions, could serve analysts who need more than one.

\paragraph{Limitations and future work.} AnchorMTM encodes one spatial direction; positions across it, and the full two- or three-dimensional geometry, are not shown on the axis. We plan to extend the bounds from one axis to a complete account of the fidelities a temporal map trades against one another, including position along several directions, temporal coherence and relative geometry, each with its own certificate. The default direction is estimated from the whole sequence, so streaming data need a direction chosen in advance. The order search is heuristic: it attains the width-aware bound on most planar frames but leaves a residual gap on PM$_{2.5}$ and on large volumetric trees. Finally, we have not yet measured how accurately analysts read position, motion and nesting from these maps; a controlled user study is the next step.

\section{Conclusion}\label{sec:conclusion}
We presented AnchorMTM, a merge tree map whose vertical axis is one fixed spatial direction, so that an analyst can read where features are and how they move while seeing how they nest. Keeping the merge tree and keeping location conflict, and we addressed this conflict in three steps: we formulated the construction of a spatially anchored map as one layout problem under a declared merge-error limit; we derived three lower bounds on the displacement that the tree forces, for exact topology, finite widths and bounded merge error; and we built the map from the data with a budgeted leaf-order search, optimal merge barriers and a capped placement. On five planar and two volumetric sequences, the selected orders attained the width-aware bound on 1,335 of 1,624 planar frames, most of the order error of existing maps proved not forced by the tree, and at the same exact topology AnchorMTM lowered the mean worst displacement by 73--96\%. Because the bounds depend only on the data, they tell before a map is drawn how much a spatial axis costs, where in the tree that cost arises, and how much merge-level fidelity buys it back. This turns the trade-off between topology and position into a quantity that can be measured and chosen, and it lays a foundation for temporal maps that certify several fidelities at once.
